\section{The finite-conductor boundary}
\label{sec:conductor-boundary}

The transfer theorem is uniform over a nontrivial coefficient class, but
that class is small compared with the prime fluctuation.  This section makes
the missing step quantitative.

For an interval
\begin{equation}\label{eq:I-N}
 I_N=(N,2N]\cap\Z,
\end{equation}
let $V_q(N)\subset\C^{I_N}$ be the restrictions of all $q$-periodic
functions, and let $P_{q,N}$ be orthogonal projection onto $V_q(N)$ in
counting-measure $\ell^2(I_N)$.  Put
\begin{equation}\label{eq:prime-fluctuation}
 g_N(n)=\Lambda(n)-1.
\end{equation}
The preceding paper proves the following estimate
\citep{WangTPCSeven2026}.

\begin{theorem}[Finite-conductor energy escape]
\label{thm:recalled-energy-escape}
For every fixed $\delta>0$, uniformly for $q\leq N^{1-\delta}$,
\begin{equation}\label{eq:projection-energy}
 \|P_{q,N}g_N\|_2^2
 \ll_\delta N\frac{q}{\varphi(q)}.
\end{equation}
Moreover,
\begin{equation}\label{eq:full-prime-energy}
 \|g_N\|_2^2\sim N\log N.
\end{equation}
Consequently,
\begin{equation}\label{eq:orthogonal-best-approximation}
 \inf_{v\in V_q(N)}\|g_N-v\|_2^2
 =\|g_N\|_2^2-\|P_{q,N}g_N\|_2^2,
\end{equation}
and hence
\begin{equation}\label{eq:best-periodic-approximation}
 \inf_{v\in V_q(N)}\|g_N-v\|_2^2
 =N\log N+O_\delta\!\left(N\frac{q}{\varphi(q)}\right)
 +o(N\log N).
\end{equation}
\end{theorem}

For $q\leq(\log N)^C$, the standard estimate
$q/\varphi(q)\ll\log\log(3q)$ gives
\begin{equation}\label{eq:polylog-energy-ratio}
 \frac{\|P_{q,N}g_N\|_2^2}{\|g_N\|_2^2}
 \ll_C\frac{\log\log\log N}{\log N}=o(1).
\end{equation}
The precise triple logarithm is not important; the vanishing ratio is.

\subsection{Why the prime coefficient cannot be inserted}

Let $v_q=P_{q,N}g_N$ and $\rho_q=g_N-v_q$.  This is the best linear
$q$-periodic approximation in $\ell^2(I_N)$.  The present theorem can
transport a bounded unit-supported periodic coefficient of period
$q\leq(\log x)^C$.  Even if one granted a separate extension of the
transfer theorem covering $v_q$, \eqref{eq:best-periodic-approximation}
shows that
\begin{equation}\label{eq:residual-energy}
 \|\rho_q\|_2^2=(1-o(1))N\log N.
\end{equation}
Thus the untransferred residual contains asymptotically all of the linear
prime-fluctuation energy.

\begin{proposition}[Strict non-specialization]
\label{prop:strict-nonspecialization}
Neither \cref{thm:operator-transfer} nor
\cref{cor:prime-source-periodic-factor} can be specialized, by periodic
projection alone, to $\beta(n)=\Lambda(n)-1$.  Doing so would require a new
estimate for the residual $\rho_q$ in \eqref{eq:residual-energy}.  The
finite-conductor transfer supplies no such estimate.
\end{proposition}

\begin{proof}
The coefficient $\Lambda-1$ is neither bounded nor periodic.  More
decisively, \eqref{eq:best-periodic-approximation} holds after optimizing
over the entire linear space of $q$-periodic functions, which is larger than
the unit-supported bounded class of \cref{thm:operator-transfer}.  Hence no
choice in that class makes the $\ell^2$ residual negligible.  The transfer
theorem contains no bilinear bound for the residual, so replacing it by zero
would be an additional unproved assertion.
\end{proof}

This is a boundary of the present method and coefficient class.  It is not a
proof that every nonlocal, nonlinear, dispersion, or Type~II construction
must fail.  In particular, an $L^2$ energy statement alone does not determine
the sign or cancellation of the residual prime-target form.

\subsection{The scale inherited from the rough prime window}

TPC-5 uses the squarefree rough modulus
\begin{equation}\label{eq:TPC5-Qw}
 Q(w)=\prod_{w<p\leq w^2}p,
 \qquad
 \log Q(w)=(1+o(1))w^2.
\end{equation}
To transfer an entire $Q(w)$-periodic observable through the present theorem,
one needs $Q(w)\leq(\log x)^C$.  Therefore
\begin{equation}\label{eq:w-polylog-scale}
 w^2\leq(C+o(1))\log\log x,
 \qquad
 w=O(\sqrt{\log\log x}).
\end{equation}
The local prime window may grow, but only to this polylogarithmic-conductor
scale under the direct Bombieri--Vinogradov proof.  Low exact-conductor
modes inside a larger $Q(w)$ do not change its ambient unit cutoff, as noted
in \cref{rem:combined-modulus-boundary}.

\subsection{A precise next-step ledger}

The current interface can be summarized as
\begin{equation}\label{eq:gap-ledger}
 \underbrace{\text{periodic low-conductor part}}_{\text{transferred here}}
 \quad+\quad
 \underbrace{\text{energy-bearing nonperiodic residual}}_{\text{uncontrolled at fixed }h}.
\end{equation}
Moving from \eqref{eq:gap-ledger} toward a fixed prime-pair problem requires
at least one genuinely new ingredient.  Possible stages are:

\begin{enumerate}[label=\textup{(\arabic*)}]
 \item replace the residue-class triangle inequality by an $L^2$ or
 coefficient-sensitive estimate and determine the largest rigorously
 transferable periodic complexity;
 \item construct a structured nonperiodic coefficient class whose residual
 has provably smaller bilinear influence, not merely smaller local Fourier
 data;
 \item obtain Type~II or dispersion control for the energy-bearing residual
 against $\Lambda(mn+h)$ at a prescribed $h$;
 \item only after such control, reconnect the exact divisor ledger to a
 fixed prime-pair correlation.
\end{enumerate}

Stages (2) and (3) are not supplied by classical Bombieri--Vinogradov.  They
are the mathematical gap, rather than an omitted routine estimate.  A
negative result showing that a proposed coefficient class cannot reduce
\eqref{eq:residual-energy} would also be a legitimate outcome of the next
stage; it would narrow the route without implying a universal parity
impossibility theorem.
